% Online Resource 1 for:
%   E. H. Brookes, "Exactness of an iterated divide-and-conquer method
%   for large-scale nonnegative least squares", Optimization Letters.
% Springer publishes supplementary PDFs as supplied; typeset to match
% the article (sn-jnl class, sn-mathphys-num style).
\documentclass[pdflatex,sn-mathphys-num]{sn-jnl}

\usepackage{amsmath,amssymb,amsfonts}
\usepackage{amsthm}
\usepackage{booktabs}

\raggedbottom
% vector math bullet (the class default introduces a Type 3 bitmap glyph)
\renewcommand{\labelitemi}{$\bullet$}

\newcommand{\supp}{\operatorname{supp}}
\newcommand{\Val}{\operatorname{Val}}
\newcommand{\T}{{\mathsf T}}
\newcommand{\prob}[1]{\mathrm{P}(#1)}

\hypersetup{
  pdftitle={Online Resource 1: An intuitive guide to the
            finite-termination and global-optimality proof},
  pdfauthor={Emre H. Brookes},
}

\begin{document}

\title[Online Resource 1]{Online Resource 1:\\ An intuitive guide to
the finite-termination and global-optimality proof}

\author*[1]{\fnm{Emre H.}
\sur{Brookes}\,\orcid{https://orcid.org/0000-0001-7348-6476}}%
\email{emre.brookes@umontana.edu}

\affil[1]{\orgdiv{Department of Chemistry \& Biochemistry},
\orgname{University of Montana}, \orgaddress{\city{Missoula},
\state{Montana}, \country{USA}}}

\maketitle

\noindent
Supplementary material for \emph{Exactness of an iterated
divide-and-conquer method for large-scale nonnegative least squares}
(Optimization Letters). This guide is explanatory: it restates the
article's argument as two short proofs by contradiction, in the
language of the motivating application. The article's numbered results
are authoritative and are cited by number; they are also
machine-checked in Lean (the article's Section~7).

\section*{1\quad The physical picture}

The measured data form a vector $b$. Each column $a_j$ of the matrix
$A$ is the predicted signal of one candidate physical component ---
in sedimentation velocity analysis, one candidate solute. We seek
nonnegative amplitudes $x_j$ (concentrations cannot be negative) so
that $Ax\approx b$, and we measure misfit by
\begin{equation*}
F(x)=\tfrac12\lVert Ax-b\rVert^2 .
\end{equation*}
The \emph{residual} $r=b-Ax$ is the part of the experiment not yet
explained by the current fit, and the \emph{support} of $x$ is the set
of components with positive amplitude.

Because the full matrix may not fit on one computer --- and because,
when few components carry signal, many small solves beat one enormous
one (the SC'06 paper quantifies this with an analytic performance
model) --- the divide-and-conquer method of the article
(Algorithms~1--2) proceeds in rounds: divide the columns into blocks; solve NNLS separately in each
block; keep each block's nonzero components; merge the survivors and
solve again; then feed the merged support back into \emph{every} block
of the next round, and repeat until nothing changes.

\section*{2\quad Comparing two answers}

Take the current fit $x$, with residual $r$, and any other candidate
answer $y$. Expanding the square,
\begin{equation*}
F(y)\;=\;F(x)\;-\;r^{\T}A(y-x)\;+\;\tfrac12\lVert A(y-x)\rVert^2 .
\end{equation*}
The last term is never negative, so $y$ can fit better than $x$ only
if $r^{\T}A(y-x)>0$: the change $y$ makes to the fitted signal must
point partly toward the unexplained data. For a single component,
$y=x+t\,e_j$ with a small $t>0$, this reads
\begin{equation*}
F(x+t\,e_j)\;=\;F(x)\;-\;t\,a_j^{\T}r\;+\;\tfrac{t^2}{2}\,
\lVert a_j\rVert^2 ,
\end{equation*}
and the sign of $a_j^{\T}r$ decides everything:
\begin{itemize}
\item if $a_j^{\T}r>0$, a small amount of component $j$ strictly
  reduces the misfit --- geometrically, $a_j$ points partly toward the
  unexplained data;
\item if $a_j^{\T}r\le0$, increasing component $j$'s amplitude cannot
  help;
\item a component already in use sits at its best amplitude exactly
  when $a_j^{\T}r=0$.
\end{itemize}
For the iterates of the algorithm the last condition is automatic:
every iterate is an \emph{exact} optimum of its restricted problem, so
each component in use already satisfies $a_j^{\T}r=0$ (Lemma~2 of the
article). For such a point, global optimality reduces to the
\emph{pricing test}: $a_j^{\T}r\le0$ for every column of the full
matrix (Lemma~3). One caution: for arbitrary $x$ the pricing test
alone is \emph{not} sufficient ($A=I$, $b=\mathbf 1$,
$x=2\cdot\mathbf 1$ has every $a_j^{\T}r<0$ yet is far from optimal),
which is why Proposition~2 of the article tracks a complementarity
residual as well.

\section*{3\quad The proof by contradiction}

Suppose the algorithm fails to find the best fit. It can fail in two
ways --- by stopping at a wrong answer, or by never stopping --- and
each leads to a contradiction. One fact serves both halves.

\paragraph{No round can make the fit worse.}
Every block of a new round contains the support of the previous merged
solution, so the previous solution is a legal candidate in every block
problem; an exactly solved block is at least as good as any legal
candidate. During merging, each child's \emph{complete} returned
support is passed upward, so each child solution remains legal at its
parent, and an exactly solved parent is at least as good as its
children. Chaining these inequalities from the leaves to the root
(Lemmas~5--6 of the article) shows that a round's merged result is at
least as good as the best of its blocks, and hence that
$v_k\le v_{k-1}$: the misfit never increases.

\paragraph{Part 1: if the algorithm stops, its answer is the best fit.}
Suppose a round leaves the result unchanged, and let $\hat x$, with
residual $\hat r$, be the answer that went into it --- the previous
round's merged result. Suppose also that some $x^*$ fits better:
$F(x^*)<F(\hat x)$.

\emph{The unexplained data point toward the better answer.} By
Section~2, $\hat r^{\T}A(x^*-\hat x)>0$. Every component that $\hat x$
uses has $a_j^{\T}\hat r=0$, so $\hat r^{\T}A\hat x=0$, and therefore
\begin{equation*}
\hat r^{\T}Ax^*\;=\;\sum_j x^*_j\,a_j^{\T}\hat r\;>\;0 .
\end{equation*}
The amplitudes $x^*_j$ are nonnegative, so at least one component $j$
that the better answer uses has $a_j^{\T}\hat r>0$.

\emph{That component gets a fair trial.} In that round --- the one
$\hat x$ went into --- $j$ lies in some block, because the blocks cover the whole matrix, and that
block also contains every component of $\hat x$, because the current
support is fed back into every block. Inside it, $\hat x$ plus a little
of component $j$ is a legal answer, and by Section~2 it fits strictly
better than $\hat x$; the exact block solve does at least as well.

\emph{The gain reaches the top.} The merged result is at least as good
as the best block, so the round's misfit is strictly below
$F(\hat x)$. Neither stopping test can then fire: the value has
changed, and an unchanged support would have forced an unchanged value
(Lemma~4(ii) of the article). This contradicts the assumption that the
round left the result unchanged. Hence $\hat x$ is a best fit, and so
is the round's own answer, which has the same misfit: the implemented
``results no longer vary'' test is a sound proof of global optimality,
not merely an empirical stopping rule (Theorem~1(b)). The article
reaches the same conclusion through the pricing test of Lemmas~3
and~7, which finds an improving column without reference to $x^*$.

\paragraph{Part 2: the algorithm does stop.}
Suppose it runs forever. The misfit never increases, and it never
stays the same, because the value test would then stop the algorithm;
so it decreases strictly, round after round. But each round's answer
is the exact optimum over some subset of the $n$ columns, and there
are at most $2^{\,n}$ subsets, hence finitely many possible misfit
values. A strictly decreasing sequence cannot run forever in a finite
set --- a contradiction. The algorithm therefore stops after finitely
many rounds, and by Part~1 it stops at a best fit (Theorem~1(a)
and~(c)). The $2^{\,n}$ bound is extremely loose; the article's
experiments stop within 2--13 rounds.

\paragraph{Why the value test matters.}
Part~1 works for either stopping test, but Part~2 needs the value
test: with the support test alone, running forever does not force the
misfit to keep falling. Under non-unique optima the misfit can sit at
its minimum while consecutive rounds keep returning different optimal
supports; the article's duplicated-column probe found such a cycle,
of period three. The support test is thus sound but not complete, and
it is the value test that guarantees halting, at round $K{+}1$
exactly, $K$ being the first optimal round (Theorem~1(c)).
Proposition~1 of the article shows why the feedback matters: a single
pass without it can lock in a strictly suboptimal answer.

\section*{4\quad What ``identical to the full solution'' means}

At termination the divide-and-conquer result has exactly the same
minimum misfit, the same fitted signal $Ax$, and the same residual
$b-Ax$ as a direct full-matrix NNLS computation, whether or not the
solution is unique (Theorem~2). The amplitude vector itself need not be unique when
columns represent identical or redundant signals; if the relevant
columns are linearly independent, the amplitudes are unique as well.

\section*{5\quad The four working assumptions}

The argument uses four properties of the implementation (conditions
(H1)--(H3) of the article, with (H2) listed here as its two clauses,
coverage and feedback):

\begin{enumerate}
\item \textbf{Exact block solves.} Every restricted NNLS problem is
  solved exactly. Without it, the inequalities of Section~3 fail; the
  article's Proposition~2 then bounds, after the fact, how far a
  returned answer can be from optimal.
\item \textbf{Full coverage.} Every round's blocks cover every column
  of the fixed full matrix. Without it, the improving component of
  Part~1 might never get its trial.
\item \textbf{Feedback of the incumbent.} The previous solution's
  complete support is included in every block. Without it, a block
  may be unable to reproduce the current solution, and both the
  no-worse step of Section~3 and Part~1 fail.
\item \textbf{Whole-support merging.} Merge stages pass complete
  returned supports upward, never splitting or truncating them
  (condition (H3) of the article). Without it, a gain found in a
  block can be lost on the way to the root, and the algorithm can
  then stop at a wrong answer (Section~8 of the article).
\end{enumerate}

\noindent In one sentence: \emph{a wrong answer cannot survive a
round, because some component of any better answer points toward the
unexplained data and gets a fair trial alongside the whole current
answer; and since only finitely many restricted optimal values exist,
the improvement must end --- at the global optimum.}

\end{document}
